% §5 Deployment Guidance --- Target 1.5pp

\section{Deployment Guidance}
\label{sec:deployment}

The per-class analysis of Sections~\ref{sec:imposs:class_fee}--\ref{sec:imposs:class_mode}
and the empirical thresholds of Section~\ref{sec:empirical} yield a
deployment matrix keyed on four parameters: (i) fee regime, (ii)
whether the watchtower batches recovery, (iii) whether the deployment
prioritises hot-key theft resistance or key-loss resilience, and (iv)
whether the deployment tolerates a federated sidechain for the
dominant-corner properties. All recommendations are derived as
counter-factual readings of the class rationality blocks
(Appendix~\ref{app:rationality}); none are asserted outside that
structure. All recommendations assume the Kerckhoffs threat model of
Section~\ref{sec:bg:lifecycle}: the attacker is taken to know the
deployment's axis values.

\begin{table}[t]
\centering
\caption{Deployment recommendations keyed on fee regime, watchtower
batching capability, and the deployment's priority (hot-key theft
resistance vs.\ key-loss resilience). Each recommendation is a
counter-factual reading of the class rationality blocks. ``Defer
to corner'' = none of the Bitcoin proposals is
class-free in this cell; operator must pick by risk weighting, or
accept federated trust to deploy on Simplicity
(Appendix~\ref{app:simplicity}).}
\label{tab:deployment}
\footnotesize
\setlength{\tabcolsep}{2.5pt}
\renewcommand{\arraystretch}{1.15}
\begin{tabular}{@{}lllll@{}}
\toprule
Fee regime                           & Batching        & Priority              & Recommended    & Rationale (class counter-factual)                                                   \\
\midrule
Low ($\leq 10\satvB{}$)              & either          & theft-resistance      & \catcsfs{}     & \classhot{} immune by \axwdbind{} dual-bind; \classscale{} N/A at low $r$           \\
Low ($\leq 10\satvB{}$)              & either          & key-loss resilience   & \opccv{}       & \classgrief{} low-cost at low $r$; key-loss covered by \axauth{} = keyless          \\
Mid ($10$--$100\satvB{}$)            & batched         & either                & \opvault{}     & \classscale{} mitigated by batching (\classscale{} §counter-factual)                \\
Mid ($10$--$100\satvB{}$)            & unbatched       & theft-resistance      & \catcsfs{}     & As above; unbatched defenders cannot safely absorb \classscale{} above $r^\dagger$  \\
Mid ($10$--$100\satvB{}$)            & unbatched       & key-loss resilience   & Defer to corner & No axis-A1 = keyless design without \classscale{} susceptibility among Bitcoin proposals \\
High ($\geq 100\satvB{}$)            & batched         & either                & \opvault{}     & Batched $r^\dagger \approx 159\satvB{}$; no anchor to pin                           \\
High ($\geq 100\satvB{}$)            & unbatched       & either                & \opctv{}       & \classfee{} cost $r$-invariant; \classscale{} N/A by \axrevault{} = no              \\
High, post-TRUC~\cite{TRUC}          & either          & either                & \opctv{}       & \classfee{} eliminated at relay-policy level; no other class active                 \\
Any                                  & any             & all classes zero      & Simplicity     & Appendix~\ref{app:simplicity}; accepts federated trust                              \\
\bottomrule
\end{tabular}
\end{table}

\paragraph{Reading the table.} Each row is a deployment scenario; the
Recommended column cites the design whose class exposures zero-out
most efficiently under the row's conditions; the Rationale column
names the specific counter-factual from the corresponding class
rationality block. Class \classfee{} (TM1) is eliminated \emph{at the
relay-policy level} by TRUC/BIP-431~\cite{TRUC}: when a \opctv{}
unvault uses a TRUC transaction type, the descendant limit collapses
to $1$ and the pinning surface vanishes. This is the single most
impactful deployment decision an \opctv{} operator can make.

\paragraph{Watchtower reserve sizing.}
For a $V = 0.5$~BTC vault at a fee rate of $100\satvB{}$ defending
against roughly one thousand unbatched \opccv{} splits, the recovery
budget is about $12.2$M~sats (${\sim}24\%$ of the vault value); the
same scenario under batching costs about $6.8$M~sats
(${\sim}14\%$), a direct consequence of the $1.8\times$
improvement in the single-block-feasibility threshold. Operators
should size reserves against the 95th-percentile fee rate of the
target deployment window and test the emergency-recovery path
quarterly.

\paragraph{Recovery-key storage policy.}
Proposition~\ref{thm:griefing_safety} forces a binary choice on axis
\axauth{}. Operators who prioritise fund safety under key loss
(\axauth{} = \texttt{keyless}, \opccv{}) accept griefing exposure
(class \classgrief{}) and must budget for the recovery-griefing
scenario (Appendix~\ref{app:rationality}, \classgrief{}). Operators
who prioritise griefing resistance (\axauth{} $\neq$
\texttt{keyless}, \opctv{}/\opvault{}/\catcsfs{}) reintroduce
key-loss risk and must use hardware-grade cold-key protection ---
particularly acute for \catcsfs{}, whose recovery leaf is
output-unconstrained (\axrecbind{} = \texttt{plain-signature},
class \classcold{}).

\paragraph{What the table does not say.}
Numerical thresholds ($r^\dagger$, reserve percentages) are
\emph{derived}, not \emph{axiomatic}. They depend on $V$, on the
watchtower's block budget, and on measured transaction sizes
(Table~\ref{tab:vsizes}). An operator with a different vault value
or a different blocks-per-recovery budget must resubstitute into
Eq.~\ref{eq:r-dagger} rather than treat this table's thresholds as
constants. The class derivations
(Sections~\ref{sec:imposs:class_fee}--\ref{sec:imposs:class_mode})
are the operator-independent content; the table is the
operator-parameterised specialisation for the reference
$V = 0.5$~BTC case.

\providecommand{\todo}[2][]{\textbf{[TODO: #2]}}
